% ============================================================================
% SUPPLEMENTARY MATERIALS: DDI RISK ANALYSIS APPLICATION
% ============================================================================

\documentclass[11pt,a4paper]{article}
\usepackage[utf8]{inputenc}
\usepackage{amsmath,amssymb}
\usepackage{booktabs}
\usepackage{geometry}
\usepackage{graphicx}
\usepackage{hyperref}
\usepackage{float}
\usepackage{longtable}
\usepackage{xcolor}
\usepackage{listings}
\geometry{margin=1in}

\hypersetup{
    colorlinks=true,
    linkcolor=blue,
    filecolor=magenta,
    urlcolor=cyan,
}

\lstset{
    basicstyle=\ttfamily\small,
    breaklines=true,
    frame=single,
    backgroundcolor=\color{gray!10}
}

\title{\textbf{Supplementary Materials}\\
\large DDI Risk Analysis Application with Alternative Recommendations}
\author{Drug-Drug Interaction Knowledge Graph Project}
\date{}

\begin{document}

\maketitle

\tableofcontents
\newpage

% ============================================================================
\section{Knowledge Graph Statistics}
% ============================================================================

\subsection{Drug Database Coverage}

The application's knowledge graph integrates multiple authoritative sources:

\begin{table}[H]
\centering
\caption{Knowledge Graph Data Sources}
\label{tab:sources}
\begin{tabular}{lrr}
\toprule
\textbf{Source} & \textbf{Interactions} & \textbf{Coverage} \\
\midrule
DrugBank 5.1.10 & 2,658,320 & Primary DDI reference \\
DDInter & 192,284 & Extended interaction pairs \\
FAERS & 149,486 & Post-market surveillance \\
\midrule
\textbf{Total (deduplicated)} & \textbf{759,774} & Multi-source validated \\
\bottomrule
\end{tabular}
\end{table}

\subsection{Drug Entity Statistics}

\begin{table}[H]
\centering
\caption{Drug Entity Coverage}
\label{tab:drugs}
\begin{tabular}{lr}
\toprule
\textbf{Metric} & \textbf{Count} \\
\midrule
Total drugs & 4,313 \\
Drugs with interactions & 5,982 \\
Drugs with side effects & 4,856 \\
Drugs with protein targets & 4,123 \\
Brand name aliases & 63 \\
\bottomrule
\end{tabular}
\end{table}

\subsection{Interaction Severity Distribution}

\begin{table}[H]
\centering
\caption{DDI Severity Classification Distribution}
\label{tab:severity_dist}
\begin{tabular}{lrr}
\toprule
\textbf{Severity} & \textbf{Count} & \textbf{Percentage} \\
\midrule
Moderate & 608,742 & 80.1\% \\
Minor & 70,682 & 9.3\% \\
Contraindicated & 44,306 & 5.8\% \\
Major & 36,044 & 4.7\% \\
\midrule
\textbf{Total} & \textbf{759,774} & 100\% \\
\bottomrule
\end{tabular}
\end{table}

% ============================================================================
\section{Drug Name Alias Dictionary}
% ============================================================================

\subsection{Brand Name Mappings}

The application maintains a curated dictionary of common brand-to-generic drug mappings to improve drug identification accuracy:

\begin{longtable}{ll}
\caption{Brand Name to Generic Drug Mappings} \\
\toprule
\textbf{Brand Name} & \textbf{Generic Name} \\
\midrule
\endfirsthead
\multicolumn{2}{c}{{\tablename\ \thetable{} -- continued}} \\
\toprule
\textbf{Brand Name} & \textbf{Generic Name} \\
\midrule
\endhead
\bottomrule
\endfoot
\bottomrule
\endlastfoot
% Common medications
Aspirin & Acetylsalicylic acid \\
Tylenol & Acetaminophen \\
Advil / Motrin & Ibuprofen \\
Coumadin & Warfarin \\
\midrule
% Statins
Lipitor & Atorvastatin \\
Zocor & Simvastatin \\
Crestor & Rosuvastatin \\
Pravachol & Pravastatin \\
Lescol & Fluvastatin \\
\midrule
% Anticoagulants
Plavix & Clopidogrel \\
Eliquis & Apixaban \\
Xarelto & Rivaroxaban \\
Pradaxa & Dabigatran \\
\midrule
% ACE Inhibitors
Cardace / Altace / Tritace & Ramipril \\
Vasotec & Enalapril \\
Zestril / Prinivil & Lisinopril \\
Lotensin & Benazepril \\
Capoten & Captopril \\
\midrule
% Beta Blockers
Lopressor / Toprol & Metoprolol \\
Tenormin & Atenolol \\
Coreg & Carvedilol \\
Inderal & Propranolol \\
\midrule
% Calcium Channel Blockers
Norvasc & Amlodipine \\
Cardizem & Diltiazem \\
Calan & Verapamil \\
\midrule
% Diuretics
Lasix & Furosemide \\
Bumex & Bumetanide \\
Demadex & Torsemide \\
\midrule
% Proton Pump Inhibitors
Nexium & Esomeprazole \\
Prilosec & Omeprazole \\
Prevacid & Lansoprazole \\
Pepcid & Famotidine \\
Zantac & Ranitidine \\
\midrule
% Antidepressants
Zoloft & Sertraline \\
Prozac & Fluoxetine \\
\midrule
% Anxiolytics
Xanax & Alprazolam \\
Valium & Diazepam \\
\midrule
% Diabetes
Glucophage & Metformin \\
Januvia & Sitagliptin \\
Jardiance & Empagliflozin \\
\midrule
% Pain/NSAIDs
Celebrex & Celecoxib \\
Voltaren & Diclofenac \\
Aleve & Naproxen \\
Ultram & Tramadol \\
Vicodin & Hydrocodone \\
Percocet & Oxycodone \\
\midrule
% Antibiotics
Augmentin & Amoxicillin \\
Zithromax & Azithromycin \\
Cipro & Ciprofloxacin \\
\end{longtable}

% ============================================================================
\section{Risk Calculation Examples}
% ============================================================================

\subsection{Sample Calculation: Warfarin + Aspirin + Ibuprofen}

Consider a regimen of warfarin, aspirin, and ibuprofen:

\subsubsection{Pairwise Interactions}

\begin{table}[H]
\centering
\caption{Interaction Analysis for Sample Regimen}
\label{tab:example}
\begin{tabular}{llr}
\toprule
\textbf{Drug Pair} & \textbf{Severity} & \textbf{Weight} \\
\midrule
Warfarin + Aspirin & Major & 0.7 \\
Warfarin + Ibuprofen & Major & 0.7 \\
Aspirin + Ibuprofen & Moderate & 0.4 \\
\midrule
\textbf{Total} & & 1.8 \\
\bottomrule
\end{tabular}
\end{table}

\subsubsection{Risk Score Calculation}

\begin{equation}
\text{Risk} = \frac{\sum w_i}{\text{max\_possible}} = \frac{1.8}{3.0} = 0.60 \quad (\text{HIGH})
\end{equation}

Where max\_possible = 3 pairs × 1.0 (contraindicated weight) = 3.0

\subsubsection{Individual PRI Scores}

\begin{table}[H]
\centering
\caption{PRI Scores for Sample Regimen}
\label{tab:pri_example}
\begin{tabular}{lrrr}
\toprule
\textbf{Drug} & \textbf{PRI} & \textbf{Risk Level} & \textbf{Alternatives Suggested} \\
\midrule
Warfarin & 0.72 & High & Yes \\
Aspirin & 0.58 & High & Yes \\
Ibuprofen & 0.41 & Medium & Yes \\
\bottomrule
\end{tabular}
\end{table}

\subsection{Alternative Recommendation Example}

For warfarin (high PRI = 0.72), the system suggests alternatives based on ARS:

\begin{table}[H]
\centering
\caption{Alternative Recommendations for Warfarin}
\label{tab:alt_example}
\begin{tabular}{llrp{5cm}}
\toprule
\textbf{Alternative} & \textbf{Class} & \textbf{ARS} & \textbf{Notes} \\
\midrule
Apixaban & DOAC & 0.78 & Lower bleeding risk, no INR monitoring \\
Rivaroxaban & DOAC & 0.71 & Once-daily dosing option \\
Dabigatran & DOAC & 0.65 & Reversible with idarucizumab \\
\bottomrule
\end{tabular}
\end{table}

% ============================================================================
\section{Application Screenshots}
% ============================================================================

\subsection{User Interface Components}

The DDI Risk Analysis Application features a three-column layout:

\begin{figure}[H]
\centering
\fbox{\parbox{0.9\textwidth}{
\textbf{[Screenshot Placeholder: Main Application Interface]}\\
\\
Three-column layout showing:
\begin{itemize}
    \item Left: Drug input methods (text, narrative, image)
    \item Center: Analysis report with risk scores and interactions
    \item Right: LLM chat interface for clinical questions
\end{itemize}
}}
\caption{DDI Risk Analysis Application Main Interface}
\label{fig:main}
\end{figure}

\begin{figure}[H]
\centering
\fbox{\parbox{0.9\textwidth}{
\textbf{[Screenshot Placeholder: Drug Validation Panel]}\\
\\
Shows identified drugs with confidence scores and edit capabilities.
}}
\caption{Drug Identification and Validation Panel}
\label{fig:validation}
\end{figure}

\begin{figure}[H]
\centering
\fbox{\parbox{0.9\textwidth}{
\textbf{[Screenshot Placeholder: Risk Analysis Report]}\\
\\
Detailed interaction report showing:
\begin{itemize}
    \item Overall polypharmacy risk score
    \item Color-coded interaction severity
    \item PRI scores for each drug
    \item Alternative recommendations with ARS
\end{itemize}
}}
\caption{Comprehensive Risk Analysis Report}
\label{fig:report}
\end{figure}

\begin{figure}[H]
\centering
\fbox{\parbox{0.9\textwidth}{
\textbf{[Screenshot Placeholder: Alternative Selection Panel]}\\
\\
Interactive panel for selecting alternative drugs and re-running analysis.
}}
\caption{Alternative Drug Selection Interface}
\label{fig:alternatives}
\end{figure}

% ============================================================================
\section{API Reference}
% ============================================================================

\subsection{Core Functions}

\begin{lstlisting}[language=Python, caption=Key Application Functions]
# Drug identification
def identify_drugs(input_text) -> (found_drugs, suggestions, not_found)

# Risk analysis
def analyze_ddi(drug_input) -> report, selection_panel, status

# PRI calculation
def calculate_regimen_pri(drug_ids) -> {drug_pris, regimen_pri}

# Alternative finding
def find_alternatives_with_ars(drug_id, other_ids) -> alternatives_list

# Chat interface
def chat(message, history, model) -> response
\end{lstlisting}

\subsection{Data Structures}

\begin{lstlisting}[language=Python, caption=Key Data Structures]
# Drug data structure
drug = {
    'drugbank_id': 'DB00001',
    'name': 'lepirudin',
    'atc_codes': ['B01AE02'],
    'description': '...'
}

# Interaction data structure
interaction = {
    'drug1_id': 'DB00001',
    'drug2_id': 'DB00002',
    'drug1_name': 'warfarin',
    'drug2_name': 'aspirin',
    'severity': 'major',
    'description': '...'
}

# PRI result structure
regimen_pri = {
    'drug_pris': {
        'warfarin': {'pri': 0.72, 'risk_level': 'High'},
        'aspirin': {'pri': 0.58, 'risk_level': 'High'}
    },
    'regimen_pri': 0.65
}
\end{lstlisting}

% ============================================================================
\section{Deployment Instructions}
% ============================================================================

\subsection{Prerequisites}

\begin{lstlisting}[language=bash, caption=Required Dependencies]
# Python 3.9+
pip install gradio pandas networkx pillow pytesseract

# LLM (LLaMA 7B via Ollama)
# Install Ollama from https://ollama.ai
# Pull LLaMA 7B model: ollama pull llama3
\end{lstlisting}

\subsection{Running the Application}

\begin{lstlisting}[language=bash, caption=Application Launch]
# Activate virtual environment
source .venv/bin/activate

# Run the application
python ddi_app.py

# Access at http://localhost:7860
\end{lstlisting}

\subsection{Configuration Options}

\begin{lstlisting}[language=Python, caption=Launch Configuration]
# Default configuration
app.launch(
    server_name="0.0.0.0",  # Listen on all interfaces
    server_port=7860,        # Default Gradio port
    share=False              # Set True for public URL
)
\end{lstlisting}

% ============================================================================
\section{Limitations and Future Work}
% ============================================================================

\subsection{Current Limitations}

\begin{enumerate}
    \item \textbf{Drug Coverage}: Limited to drugs present in DrugBank database
    \item \textbf{Interaction Evidence}: Severity classifications based on database labels, not real-time evidence
    \item \textbf{OCR Accuracy}: Image extraction depends on prescription legibility
    \item \textbf{LLM Dependency}: Chat functionality requires API access or local model setup
\end{enumerate}

\subsection{Planned Enhancements}

\begin{enumerate}
    \item Integration with additional drug databases (RxNorm, SNOMED-CT)
    \item Real-time FAERS signal monitoring
    \item Patient-specific risk adjustment (age, renal function, genetics)
    \item Mobile-responsive interface for bedside use
    \item Electronic health record (EHR) integration via FHIR
\end{enumerate}

\end{document}
